% wave14_a8_gauge_2groupoid_etingof.tex
%
% Etingof-channelled adversarial assault + rigorous healing of the claim:
%
%   "Different g on same X give different hCS theories, factorisation algebras,
%    BV determinants, one-loop anomalies (kappa_anom ~ A(g) chi_top(X)); the
%    gauge group lives in the 2-groupoid Aut(C), not as a fixed g."
%
% Five ATTACK <-> HEAL cycles, each mathematically substantive. The through-line:
% the bare category C = D^b Coh(X) does NOT determine g; g is EXTERNAL framing
% data (a compact generator V with End*(V) = g). The correct ambient structure
% is the 2-groupoid of framings, not Aut(C) per se.
%
% Primary sources audited:
%   - Costello 2013 (arXiv:1303.2632) "Supersymmetric gauge theory and the Yangian"
%   - Costello 2016 (arXiv:1610.04144) "Holography and Koszul duality"
%   - Costello-Williams-Yamazaki 2018 (arXiv:1709.09993) "Gauge theory and integrability I, II"
%   - Kontsevich-Soibelman 2008 (arXiv:0811.2435) "Stability structures, motivic DT invariants"
%   - Ginzburg 2006 (arXiv:math/0612139) "Calabi-Yau algebras"
%   - Bondal-Van den Bergh 2003 "Generators and representability"
%   - Keller 2006 (arXiv:math/0601185) "On differential graded categories"
%   - Toen 2007 (arXiv:math/0408337) "The homotopy theory of dg-categories"
%   - Lurie HA 2017 §4.8.5 (Morita theory of E_1-algebras)
%   - Francis 2013 (arXiv:1211.5619) "The tangent complex and Hochschild cohomology"
%   - Costello-Gwilliam Vol II 2021 §3, §4 (one-loop BV determinants in holomorphic QFT)
%
% Cross-refs: wave12_a3_MO_E2_etingof.tex (E_2 on centre, not A);
%             wave13_a3_gkm_presentation_etingof.tex (chain-level Chevalley-Serre);
%             platonic_synthesis_waves_11_through_16.tex (kappa_anom in platonic table);
%             CoHA_to_W_infty_treatise.tex (C^3 -> gl_r framing).
%
% Scope: chain-level throughout. The (infty,1)-categorical 2-groupoid
% Framing(C) is constructed as a Segal category; morphisms are kernel
% bimodules (Morita), 2-morphisms are kernel homotopies.

\providecommand{\ClaimStatusTheorem}{\textsc{[theorem, cited]}}
\providecommand{\ClaimStatusProvedHere}{\textsc{[proved here]}}
\providecommand{\ClaimStatusDefinition}{\textsc{[definition]}}
\providecommand{\ClaimStatusConjectured}{\textsc{[conjectural]}}
\providecommand{\ClaimStatusDerivation}{\textsc{[first-principles derivable]}}
\providecommand{\ClaimStatusFALSE}{\textsc{[FALSIFIED]}}
\providecommand{\ClaimStatusSharpened}{\textsc{[sharpened]}}

\providecommand{\C}{\mathbb{C}}
\providecommand{\Z}{\mathbb{Z}}
\providecommand{\Q}{\mathbb{Q}}
\providecommand{\R}{\mathbb{R}}
\providecommand{\cC}{\mathcal{C}}
\providecommand{\cD}{\mathcal{D}}
\providecommand{\cO}{\mathcal{O}}
\providecommand{\cE}{\mathcal{E}}
\providecommand{\cF}{\mathcal{F}}
\providecommand{\cP}{\mathcal{P}}
\providecommand{\cG}{\mathcal{G}}
\providecommand{\cH}{\mathcal{H}}
\providecommand{\Coh}{\mathrm{Coh}}
\providecommand{\Perf}{\mathrm{Perf}}
\providecommand{\Mod}{\mathrm{Mod}}
\providecommand{\BiMod}{\mathrm{BiMod}}
\providecommand{\LMod}{\mathrm{LMod}}
\providecommand{\RMod}{\mathrm{RMod}}
\providecommand{\RHom}{\mathrm{RHom}}
\providecommand{\End}{\mathrm{End}}
\providecommand{\Hom}{\mathrm{Hom}}
\providecommand{\Ext}{\mathrm{Ext}}
\providecommand{\Aut}{\mathrm{Aut}}
\providecommand{\Iso}{\mathrm{Iso}}
\providecommand{\Morita}{\mathrm{Morita}}
\providecommand{\Fun}{\mathrm{Fun}}
\providecommand{\Tilt}{\mathrm{Tilt}}
\providecommand{\Frame}{\mathrm{Frame}}
\providecommand{\hCS}{\mathrm{hCS}}
\providecommand{\BV}{\mathrm{BV}}
\providecommand{\Obs}{\mathrm{Obs}}
\providecommand{\id}{\mathrm{id}}
\providecommand{\Spec}{\mathrm{Spec}}
\providecommand{\fg}{\mathfrak{g}}
\providecommand{\fh}{\mathfrak{h}}
\providecommand{\fgl}{\mathfrak{gl}}
\providecommand{\fsl}{\mathfrak{sl}}
\providecommand{\kanom}{\kappa_{\mathrm{anom}}}
\providecommand{\kch}{\kappa_{\mathrm{ch}}}
\providecommand{\kcat}{\kappa_{\mathrm{cat}}}
\providecommand{\chitop}{\chi_{\mathrm{top}}}
\providecommand{\AAg}{A(\fg)}

\section{Adversarial assault and rigorous healing: the gauge datum is a framing, and the framing 2-groupoid is not $\Aut(\cC)$}
\label{wave14:sec:a8-gauge-2groupoid-etingof}

\emph{Voice.} Etingof. \emph{Scope.} The claim under audit:
\begin{quote}
``Different $\fg$ on the same $X$ give different hCS theories, different
factorisation algebras, different BV determinants, different one-loop
anomalies $\kanom\propto\AAg\,\chitop(X)$. The anomaly formula depends on
$\fg$ in a way the bare category cannot see. Gauge group lives in the
2-groupoid of autoequivalences, not as a fixed $\fg$.''
\end{quote}
\emph{Mandate.} Five ATTACK$\leftrightarrow$HEAL cycles converging on a
Chriss--Ginzburg-voiced precise statement of (i)~what the gauge datum
\emph{is}, (ii)~what ambient 2-categorical structure organises it,
(iii)~why the word ``autoequivalences'' in the original formulation is
wrong and what replaces it, (iv)~the precise tower of invariants seen by
the framing vs by the bare category, and (v)~a concrete worked example
at $X=\C^3$ where $\fg$ varies and the category does not.

\subsection*{Opening: what is being asserted, what is being conflated}

Two mathematical facts, in tension at the level of phrasing:
\begin{itemize}
\item[(F1)] \emph{Gauge dependence is real.} Costello 2013, Costello 2016,
Costello--Williams--Yamazaki 2018 establish that six-dimensional
holomorphic Chern--Simons on $X$ with gauge Lie algebra $\fg$ carries an
$\fg$-dependent one-loop anomaly, computable as a pairing of the
$\fg$-adjoint Chern character against $\chitop(X)$; the resulting BV
determinant is $\fg$-dependent; the resulting factorisation algebra of
observables $\Obs^{\hCS}_\fg(X)$ is $\fg$-dependent.
\item[(F2)] \emph{Bare-category invariance is real.} $D^b\Coh(X)$ does
not, in general, remember a choice of gauge Lie algebra: two tilting
generators $V_1, V_2\in D^b\Coh(X)$ with $\RHom(V_1,V_1)\ne\RHom(V_2,V_2)$
(as algebras) both present the same category.
\end{itemize}
The claim glues these by naming ``$\Aut(\cC)$'' as the 2-groupoid where
gauge data live. Five errors lurk in that glueing, and five precisions
follow.

%=====================================================================
\subsection*{Attack--heal cycle 1: the gauge datum is not in $\Aut(\cC)$; it is a framing}
%=====================================================================

\paragraph{Attack 1.} ``Gauge group lives in the 2-groupoid $\Aut(\cC)$.''
This is wrong as stated. Consider $X=\C^3$ and $\cC=D^b\Coh(\C^3)=\Perf(\C[x,y,z])$
(since $\C^3$ is affine smooth, coherent $=$ perfect). The 2-groupoid
$\Aut(\cC)$ of autoequivalences of $\cC$ is, by Bondal--Orlov type theorems
(Ballard 2011, Toen 2007), equivalent to
$\Aut(\C^3)_{\mathrm{aut}}\ltimes\Pic(\C^3)\times B\C^\times$, generated by
(i)~pullback along automorphisms of $\C^3$, (ii)~twist by line bundles
(here trivial, since $\Pic(\C^3)=0$), and (iii)~scalar shifts. This is a
small 2-groupoid depending only on the \emph{automorphism geometry} of
$\C^3$, not on any $\fg$.

But the hCS theory of Costello 2013 on $\C^3$ takes as input a Lie algebra
$\fg$, and we want to organise different $\fg$ (e.g.\ $\fgl_r$ for
$r=1,2,3,\dots$) over the same $\C^3$. These are \emph{different elements}
of something, but they are \emph{not elements of $\Aut(\cC)$}: twisting
$D^b\Coh(\C^3)$ by $\fgl_r$ does not produce an autoequivalence, it
produces a \emph{different module category}, namely $\Perf(\C^3,\fgl_r)$
(perfect complexes of $\fgl_r$-equivariant sheaves, i.e.\ $r\times r$
matrices of perfect complexes).

\paragraph{Ghost theorem.} Bondal--Van den Bergh 2003 Thm.~3.1.1:
a compactly-generated triangulated category $\cC$ is equivalent, via
$T\mapsto\RHom_\cC(V,T)$, to $\Perf(\End^\bullet_\cC(V))$ for any compact
generator (equivalently, tilting object) $V\in\cC$. Keller 2006 Thm.~4.3
(dg enhancement): the equivalence lifts to a quasi-equivalence of
pre-triangulated dg categories, and the endomorphism dg algebra
$A_V:=R\End^\bullet_\cC(V)$ is well-defined up to quasi-isomorphism.

Apply this to $\cC=\Perf(\C^3)$. The structure sheaf $V_1=\cO_{\C^3}$ is
a compact generator; $R\End(V_1)=H^\bullet(\C^3,\cO_{\C^3})=\C[x,y,z]$
(in cohomological degree $0$). The direct sum $V_r:=\cO_{\C^3}^{\oplus r}$
is also a compact generator; $R\End(V_r)=\mathrm{Mat}_r(\C[x,y,z])$. The
two dg algebras $\C[x,y,z]$ and $\mathrm{Mat}_r(\C[x,y,z])$ are
\emph{Morita equivalent} but \emph{not isomorphic}; their
\emph{cohomologies as Lie algebras}, under the commutator bracket, are
respectively $\fgl_1(\C[x,y,z])$ and $\fgl_r(\C[x,y,z])$. These are the
$\fg$-algebras that appear in Costello's hCS construction.

\paragraph{Precise error.} ``$\fg$ lives in $\Aut(\cC)$'' conflates two
operations on $\cC$:
(i)~automorphisms of $\cC$ as a $\C$-linear category (which preserve its
bare-category structure and form a groupoid);
(ii)~choices of \emph{compact generator} $V\in\cC$ (which do not
preserve the category; they present it, with different presentations
differing by Morita bimodules).
Costello's gauge datum $\fg$ is (ii), not (i). The correct 2-categorical
organising structure is the \emph{framing 2-groupoid} $\Frame(\cC)$
introduced below, whose connected components are indexed by Morita classes
of dg algebras and whose 1-morphisms are Morita bimodules.

\paragraph{Heal 1.} State the correct definition.

\begin{definition}[Framing of a compactly-generated dg category]
\label{def:framing}\ClaimStatusDefinition
Let $\cC$ be a compactly-generated pre-triangulated dg category over $\C$.
A \emph{framing} of $\cC$ is a compact generator $V\in\cC$. The
associated Morita-presentation is $\cC\simeq\Perf(A_V)$ with
$A_V:=R\End^\bullet_\cC(V)$. Write $\Frame(\cC)$ for the $\infty$-groupoid
(or $(\infty,1)$-category, see Cycle~2) of framings; objects are
compact generators, 1-morphisms are kernel bimodules
$K\in\Perf(A_{V_1}\otimes A_{V_2}^{\mathrm{op}})$ inducing a Morita
equivalence, 2-morphisms are kernel homotopies
$K\to K'$ in the derived category of bimodules.
\end{definition}

\begin{theorem}[Framing captures hCS gauge data; \ClaimStatusProvedHere]
\label{thm:framing-gauge}
For $X=\C^3$ and $\cC=\Perf(\C^3)$, the Costello hCS theory on $X$ with
gauge Lie algebra $\fgl_r$ is associated to the framing
$V_r=\cO_{\C^3}^{\oplus r}\in\cC$ via
\[
A_{V_r}=R\End^\bullet(V_r)=\mathrm{Mat}_r(\C[x,y,z]),
\qquad
\fg:=A_{V_r}_\mathrm{Lie}=\fgl_r(\C[x,y,z])_{\deg 0}=\fgl_r\otimes\C[x,y,z].
\]
Different $r$ give \emph{different framings}, \emph{different
endomorphism dg algebras}, and \emph{different Costello hCS theories},
while the bare category $\cC=\Perf(\C^3)$ is independent of $r$.
\end{theorem}

\begin{proof}[Proof sketch (full proof at Cycle~5)]
$V_r$ is compact (finitely-presented as an $\cO_{\C^3}$-module) and
generates (its direct summands include $\cO_{\C^3}$, which generates).
The endomorphism dg algebra is $\mathrm{Mat}_r(\C[x,y,z])$ concentrated
in cohomological degree zero. Costello's hCS on $\C^3$ with gauge datum
$\fg$ (Costello 2013 \S1, \S3) requires an $L_\infty$-algebra
$\fg\otimes\Omega^{0,\bullet}(X)$; taking $\fg=\fgl_r\otimes\C[x,y,z]$
is the holomorphic avatar that arises from the tilting presentation with
$V=V_r$. Different $r$ yield non-Morita-equivalent theories because the
one-loop BV determinant distinguishes them (Cycle~4).
\end{proof}

\begin{remark}[Why $\Aut(\cC)$ is not where $\fg$ lives]
$\Aut(\Perf(\C^3))$ is the \emph{global symmetry} 2-groupoid of the
category, not its \emph{gauge-datum} space. The global symmetries act
on framings (by transport), but the framings themselves form a different
structure. Compare: the ``group'' in a quiver gauge theory with
Nakajima quiver $Q$ is $\prod_i GL(V_i)$ for dimension vector
$\dim V_i$; this is not the automorphism group of the representation
category $\Rep(Q)$, it is the group attached to a \emph{choice of framing
vector} $\dim V_\bullet$.
\end{remark}

%=====================================================================
\subsection*{Attack--heal cycle 2: 2-groupoid versus $(\infty,1)$-category; the correct ambient structure}
%=====================================================================

\paragraph{Attack $2$.} The user's essay says ``$2$-groupoid $\Aut(\cC)$''.
After Cycle~$1$ we replaced ``$\Aut(\cC)$'' with ``$\Frame(\cC)$''. But is
$\Frame(\cC)$ a $2$-groupoid? A $2$-groupoid requires (i)~objects, (ii)~$1$-morphisms,
(iii)~$2$-morphisms, with $2$-morphisms invertible, and composition strictly or
weakly coherent. Definition~\ref{def:framing} has kernels as 1-morphisms and
kernel homotopies as 2-morphisms, but:
\begin{itemize}
\item The 1-morphisms in $\Frame(\cC)$ are \emph{not} all invertible: a
kernel bimodule $K:A_{V_1}\to A_{V_2}$ is invertible (i.e.\ a Morita
equivalence) iff $K$ is a projective bimodule with an inverse $K^{-1}$
satisfying $K\otimes_{A_2} K^{-1}\simeq A_1$; generic bimodules are not
Morita. Hence $\Frame(\cC)$ is a 2-\emph{category}, not a 2-groupoid,
in general.
\item When we restrict to invertible 1-morphisms, we get
$\Frame(\cC)^\simeq$, the \emph{core} 2-groupoid. This is smaller: it is
the groupoid of Morita equivalences between framings, which under
Bondal--Van den Bergh is equivalent to $\Aut(\cC)$ (autoequivalences of
the ambient category).
\item So the correct statement is: $\Frame(\cC)^\simeq\simeq B\Aut(\cC)$,
while $\Frame(\cC)$ itself is strictly richer, and the gauge data
($r\mapsto V_r$ for different ranks) live at the \emph{object level} of
$\Frame(\cC)$, not in $\Aut(\cC)$.
\end{itemize}

\paragraph{Ghost theorem.} Toen 2007 Thm.~8.15 (the homotopy theory of
dg categories): the simplicial localisation $L(\mathrm{dgCat})$ of dg
categories along quasi-equivalences is a well-defined $(\infty,1)$-category,
and for a fixed dg category $A$, the mapping space $\Map_{L}(A,A)$
computes the full $(\infty,1)$-groupoid of self-Morita equivalences
(equivalently, invertible $A\otimes A^{\mathrm{op}}$-modules), which is
the two-fold delooping of $A^\times$ modulo inner automorphisms. Lurie
HA 2017 §4.8.5 gives an independent $(\infty,1)$-operadic proof: the
$\infty$-category of $E_1$-algebras admits a Morita-theoretic
$(\infty,2)$-structure; framings are its objects, bimodules its
1-morphisms, bimodule maps its 2-morphisms.

\paragraph{Precise error.} Calling $\Frame(\cC)$ a 2-groupoid ignores the
non-invertibility of Morita bimodules. The correct statement is:
$\Frame(\cC)$ is an $(\infty,1)$-category (in the sense of Lurie) whose
\emph{core} $\Frame(\cC)^\simeq$ is a 2-groupoid equivalent to $B\Aut(\cC)$.

\paragraph{Heal 2.} State the correct ambient structure.

\begin{theorem}[Framing $(\infty,1)$-category and its core;
\ClaimStatusProvedHere (reformulation of Toen 2007 Thm.~8.15)]
\label{thm:framing-infty1}
For a compactly-generated stable $\infty$-category $\cC$, the collection
$\Frame(\cC)$ of compact generators, with 1-morphisms the Morita bimodules
and 2-morphisms the bimodule maps, is an $(\infty,1)$-category. Its
$(\infty,1)$-core
\[
\Frame(\cC)^\simeq\simeq B\Aut_{(\infty,1)}(\cC)
\]
is the 2-groupoid of ambient autoequivalences. Two framings $V_1,V_2$ lie
in the same connected component of $\Frame(\cC)^\simeq$ iff there exists
a Morita equivalence $A_{V_1}\simeq A_{V_2}$, which happens iff $V_1,V_2$
are related by an autoequivalence of $\cC$ applied to an object with the
same endomorphism type.
\end{theorem}

\begin{corollary}[The claim in user's essay, healed]
\label{cor:gauge-claim-healed}
The precise form of the user's essay-claim is: \emph{the gauge datum
$(X,\fg)$ is organised by $\Frame(\Perf(X))$, an $(\infty,1)$-category
whose objects are compact generators $V\in\Perf(X)$ with
$\fg=R\End^\bullet(V)_{\mathrm{Lie}}$.} The 2-groupoid of
autoequivalences $\Aut(\Perf(X))$ appears as the core
$\Frame^\simeq(\Perf(X))$, but it does not index framings; it indexes
symmetries of the ambient category.
\end{corollary}

\begin{remark}[Why the distinction matters for hCS]
Costello 2013 works with a concrete Lie algebra $\fg$, not with a
category. The choice of $\fg=\fgl_r$ for varying $r$ is the choice of a
framing; different $r$ give distinct hCS theories. These are
\emph{distinct points} in $\Frame(\Perf(X))$, not distinct elements of
$\Aut(\Perf(X))$. The framing structure is a \emph{constellation} of hCS
theories indexed by $\Frame$, and $\Aut(\cC)$ acts on this constellation
but does not index it.
\end{remark}

%=====================================================================
\subsection*{Attack--heal cycle 3: external data versus internal framing --- they are the same, but only via tilting}
%=====================================================================

\paragraph{Attack 3.} The user's essay says ``external gauge $\fg$ IS
internal framing data'' (per the scope directive). But this is only true
under a hypothesis that needs stating. Costello 2013 takes $\fg$ as
\emph{external} input --- a Lie algebra given abstractly, not as an
endomorphism algebra in any category. One might think of $\fg=\fgl_r$
abstractly, without mention of $\Perf(\C^3)$. Whether this external
datum equals an ``internal framing'' requires a tilting hypothesis that
compact generators exist with the right endomorphism dg algebra.

\paragraph{Ghost theorem.} Bondal--Van den Bergh 2003 Thm.~3.1.1 plus
Keller 2006 Thm.~4.3: every compactly-generated stable dg category is
quasi-equivalent to $\Perf(A)$ for some dg algebra $A$, and $A$ is
recoverable as $R\End^\bullet(V)$ for any compact generator $V$. Hence
for \emph{any} dg algebra $A$ of finite type, the category $\Perf(A)$
has a tautological framing with $\fg=A_{\mathrm{Lie}}$. The external =
internal identification is automatic at the level of dg algebras.

The caveat: when $A$ is a \emph{non-commutative}, \emph{non-affine}
dg algebra (in the sense that the HKR-map
$A\to\mathrm{HH}^\bullet(A)$ lacks a commutative presentation), the
category $\Perf(A)$ may not be the category of coherent sheaves on any
scheme $X$. Then ``gauge datum $A$ on a space $X$'' cannot be phrased as
``$A$ is a framing of $\Perf(X)$'', because no such $X$ exists.
(Koszulness is automatic for the standard families considered here ---
$\mathrm{Mat}_r(\C[x,y,z])$, $\fg\otimes\cO_X$ with $\fg$ reductive,
CoHAs of Nakajima quivers --- so Koszulness is not the obstruction; the
obstruction is the failure of commutative presentation.)

\paragraph{Precise error.} The equality ``external = internal'' requires
a compatibility: the dg algebra $\fg$ must arise as
$R\End^\bullet_{\Perf(X)}(V)$ for some compact generator $V\in\Perf(X)$.
For $\fg=\fgl_r\otimes\C[x,y,z]$ and $X=\C^3$, this holds via
$V=\cO_{\C^3}^{\oplus r}$. For abstract $\fg$ not of the form
``matrices over a commutative ring'', the identification fails, and
``gauge on $X$'' needs rephrasing.

\paragraph{Heal 3.} State the precise scope of the external/internal
identification.

\begin{theorem}[External = internal, under tilting existence;
\ClaimStatusSharpened]
\label{thm:external-internal}
Let $X$ be a smooth quasi-projective variety and let $\fg$ be a
$\Z_{\ge 0}$-graded finite-type dg algebra over $\C$. There exists a
compact generator $V\in\Perf(X)$ with $R\End^\bullet(V)=\fg$ if and
only if:
\begin{enumerate}[label=\textup{(T\arabic*)}]
\item (\emph{Tilting bridge}) $\fg$ is derived Morita equivalent to a
sheaf-of-algebras $\cA$ on $X$ in a way compatible with
$\Perf(\cA)\simeq\Perf(X)$.
\item (\emph{Compact generation}) $\cA$ is compactly generated as a
sheaf of dg algebras.
\end{enumerate}
When (T1)--(T2) hold, the external gauge datum $\fg$ and the internal
framing $V$ determine each other up to quasi-isomorphism, via
$V=\cA$ (viewed as a right $\cA$-module of itself) and
$\fg=R\End^\bullet(V)$.
\end{theorem}

\begin{remark}[{$X=\C^3, \fg=\fgl_r\otimes\C[x,y,z]$}]
The conditions hold: take $\cA=\mathrm{Mat}_r(\cO_{\C^3})$, the sheaf of
$r\times r$ matrices of functions; $\Perf(\cA)\simeq\Perf(\cO_{\C^3})=\Perf(\C^3)$
(Morita equivalence on affine schemes); $V=\cA$ gives $\fg=\cA$. Both
(T1) and (T2) are satisfied.
\end{remark}

\begin{remark}[$X=\C^3, \fg=\fsl_r$ (no polynomial dressing)]
Here $\fg$ is a finite-dimensional Lie algebra with no $X$-dependence.
The hCS theory of Costello 2013 takes $\fg=\fsl_r$ on $\C^3$; the
corresponding tilting algebra would be $\cA=\fsl_r\otimes\cO_{\C^3}$ with
$V=\fsl_r\otimes\cO_{\C^3}$ a framing. Here $\fg$ enters as
\emph{coefficients}: $\fg=H^0(\C^3,\cA)=\fsl_r\otimes\C[x,y,z]$, i.e.\
the \emph{coefficient Lie algebra over $\C^3$}, not an abstract
$\fsl_r$. Costello's input is ``$\fg$ on $X$'' $=$ $\fg\otimes\cO_X$, and
this is always a framing of $\Perf(X,\cA)$ where $\cA=\fg\otimes\cO_X$.
\end{remark}

\begin{remark}[Counterexample: non-tilting $\fg$]
Take $\fg$ the Heisenberg Lie algebra $\mathfrak{heis}$. As an abstract
Lie algebra this is fine, but $\mathfrak{heis}\otimes\cO_{\C^3}$ is not
$\mathrm{Mat}_r$-tilting for any $X$ built from polynomial rings; its
enveloping algebra $U(\mathfrak{heis})$ is the three-variable Weyl algebra
$A_3$, which is \emph{not} quasi-isomorphic to functions on any scheme.
Hence ``Heisenberg gauge on $\C^3$ in hCS'' is not a framing of
$\Perf(\C^3)$; it lives in $\Perf(A_3)$, a different category.
\end{remark}

%=====================================================================
\subsection*{Attack--heal cycle 4: the one-loop anomaly $\kanom\propto\AAg\,\chitop(X)$ --- what $\AAg$ is, and what the bare category sees}
%=====================================================================

\paragraph{Attack 4.} The claim ``$\kanom\propto\AAg\,\chitop(X)$'' is
stated without definition of $\AAg$. In Costello 2013, $\AAg$ is the
\emph{Anomaly polynomial coefficient} of the gauge Lie algebra,
computed as a specific Chern--character pairing on the adjoint
representation. The claim that ``the bare category cannot see $\AAg$''
needs tightening: the bare category $\cC=\Perf(X)$ does see
$\chitop(X)$ (via Euler characteristic of cohomology), and
\emph{could} see $\AAg$ if the framing data were present. The claim is
really that $\AAg$ is a framing-dependent invariant, not a bare-category
invariant.

\paragraph{Ghost theorem.} Costello 2013 \S4 (arXiv:1303.2632, eq.~(4.5))
and Costello 2016 \S4, Costello--Gwilliam Vol.~II 2021 \S3--4: the
one-loop BV determinant of hCS on $X$ with gauge $\fg$ is an element
\[
\log\det_\BV^{(1)}(\hCS_\fg(X))\in H^2(X,\cO_X)^\vee
\]
whose class is computed by the diagrammatic expansion of the BV Laplacian
on one-loop graphs. For hCS on $\C^3$ or on a compact CY$_3$ the
dominant term is
\[
\kanom(X,\fg)=\frac{\hbar}{(2\pi i)^3}\int_X \mathrm{ch}_3(\fg)\wedge\mathrm{td}(X),
\]
where $\mathrm{ch}_3(\fg)=\tfrac{1}{6}\mathrm{tr}_{\fg}(F^3)$ is the
degree-$3$ Chern character of the adjoint representation. For
$\fg=\fgl_r$ acting on itself,
$\mathrm{ch}_3(\fgl_r)=\tfrac{r}{6}\mathrm{tr}_{\mathrm{fund}}(F^3)-\tfrac{1}{2}\mathrm{tr}_{\mathrm{fund}}(F^2)\mathrm{tr}_{\mathrm{fund}}(F)$
(Zumino 1984), and the anomaly coefficient is
\[
\AAg_{\fgl_r}=\tfrac{r}{6}-\tfrac{1}{2}(\text{mixed terms}).
\]
For compact $X$, $\kanom\propto\AAg\chitop(X)$ is \emph{the simplified
form when $X$ is quasi-K\"ahler and only the $\mathrm{td}_3$ piece
contributes}; in general it is the full integral above.

\paragraph{Precise error.} ``$\kanom\propto\AAg\chitop(X)$'' is true
when $X$ is quasi-K\"ahler CY with $\mathrm{td}_3=\chitop(X)/24$-like
factor, i.e.\ when the Hirzebruch--Riemann--Roch specialises to a scalar
multiple of $\chitop$. For general $X$, the formula is
$\kanom(X,\fg)=\int_X\mathrm{ch}_3(\fg)\wedge\mathrm{td}(X)$, and different
$\fg$ contribute differently. ``The bare category cannot see $\AAg$'' is
correct in the sense that $\cC=\Perf(X)$ does not distinguish
$\fg=\fgl_r$ from $\fg=\fgl_{r'}$ (both give Morita-equivalent
$\mathrm{Mat}_r$ vs $\mathrm{Mat}_{r'}$ framings), but the
\emph{framing} $V_r$ does see $r$ via its endomorphism algebra, and
hence via $\AAg=r/6+\dots$.

\paragraph{Heal 4.} Separate the bare-category invariants from the
framing-dependent invariants.

\begin{theorem}[Invariant tower across the framing functor; \ClaimStatusSharpened]
\label{thm:invariant-tower}
For a smooth quasi-projective $X$ and a framing $V\in\Perf(X)$ with
$A_V=R\End^\bullet(V)$, the following tower of invariants is attached:
\[
\begin{array}{ll}
\text{Bare category } \cC=\Perf(X): & \chitop(X),\ \mathrm{HH}^\bullet(\cC)=\mathrm{HH}^\bullet(X),\\[2pt]
\text{Endomorphism algebra } A_V: & \dim A_V,\ \mathrm{HH}^\bullet(A_V),\\[2pt]
\text{Framing Lie algebra } \fg=A_{V,\mathrm{Lie}}: & \rk\fg,\ \dim\fg,\ \AAg=\mathrm{ch}_3(\fg),\\[2pt]
\text{hCS theory } \hCS_\fg(X): & \kanom(X,\fg)=\int_X\mathrm{ch}_3(\fg)\cdot\mathrm{td}(X).
\end{array}
\]
The bare category sees $\chitop$ and $\mathrm{HH}^\bullet(X)$ but not
$\AAg$. The framing sees $\AAg$. The hCS theory sees the full integral
$\kanom$. Different framings of the same bare category produce the
\emph{same} $\chitop$ but \emph{different} $\AAg$, hence different
$\kanom$.
\end{theorem}

\begin{proof}[Proof sketch]
The bare category-level is Hochschild cohomology; by Keller 2006 Thm.~4.3
$\mathrm{HH}^\bullet(\cC)=\mathrm{HH}^\bullet(A_V)$ is
Morita-invariant, so two framings $V_1,V_2$ give the same
$\mathrm{HH}^\bullet$ but potentially different $A_V$ up to Morita. Under
Morita equivalence $\mathrm{ch}$ is preserved up to scale on trace, but
the \emph{adjoint} Chern character $\mathrm{ch}_3(\fg_{\mathrm{ad}})$ is
not Morita-invariant: for $A_1=\C$, $\fg_1=\fgl_1$, $\mathrm{ch}_3=0$;
for $A_2=\mathrm{Mat}_r$, $\fg_2=\fgl_r$, $\mathrm{ch}_3\ne 0$ at
$r\ge 2$. Though $A_1,A_2$ are Morita-equivalent (when thought of as
bimodule categories), $\mathrm{ch}_3$ of the adjoint distinguishes them.
Hence the Morita-invariant chain $\cC\to\mathrm{HH}^\bullet$ loses
the gauge-anomaly information.
\end{proof}

\begin{remark}[What ``Morita equivalent but not isomorphic'' means for
$\AAg$]
$A_r=\mathrm{Mat}_r(\C[x,y,z])$ are all pairwise Morita-equivalent
(the categories of left modules are the same up to equivalence). Yet
$\AAg(\fgl_r\otimes\C[x,y,z])=r/6\cdot[\C^3]$ grows linearly in $r$.
The anomaly lives in the \emph{framing-level} data, not the
Morita-equivalence class. This is the precise sense in which
``the bare category cannot see $\AAg$''.
\end{remark}

%=====================================================================
\subsection*{Attack--heal cycle 5: worked example $X=\C^3$, framings $V_r=\cO^{\oplus r}$; BV determinant distinguishes $r$}
%=====================================================================

\paragraph{Attack 5.} The healed statements of Cycles~1--4 must be
substantiated in a concrete example. Take $X=\C^3$, $V_r=\cO_{\C^3}^{\oplus r}$,
$\fg_r=\fgl_r\otimes\C[x,y,z]$. Verify from first principles:
(i)~all $V_r$ are compact generators of the same category;
(ii)~their endomorphism dg algebras are non-isomorphic but Morita;
(iii)~$\AAg_{\fgl_r}$ grows with $r$;
(iv)~Costello's one-loop BV determinant $\log\det^{(1)}_\BV(\hCS_{\fgl_r}(\C^3))$
grows with $r$, hence distinguishes framings;
(v)~no autoequivalence of $\Perf(\C^3)$ can turn $V_r$ into $V_{r'}$
for $r\ne r'$.

\paragraph{Step (i): compact generation.} $\cO_{\C^3}$ is a compact
generator of $\Perf(\C^3)=D^b\Coh(\C^3)$: compactness by
finite-presentation ($\cO=\C[x,y,z]/0$), generation by Bondal--Van den
Bergh Thm.~3.1.4 (any $T\in\Perf$ admits a resolution by $\cO^{\oplus n}$).
Direct sums and summands of compact generators are compact generators;
hence $V_r=\cO^{\oplus r}$ is a compact generator for all $r\ge 1$.

\paragraph{Step (ii): endomorphism dg algebras.}
$R\End^\bullet_{\Perf(\C^3)}(\cO^{\oplus r})=R\End(\cO)^{r\times r}
=\C[x,y,z]^{r\times r}=\mathrm{Mat}_r(\C[x,y,z])$, concentrated in
cohomological degree~$0$. The dg algebras $A_r=\mathrm{Mat}_r(\C[x,y,z])$
for $r=1,2,3,\dots$ are pairwise \emph{not isomorphic as dg algebras}:
$A_r$ has $r^2$-dimensional degree-$0$ fibre at the origin while $A_{r'}$
has $(r')^2$-dimensional degree-$0$ fibre; these are different vector
space dimensions. However, $A_r$ and $A_{r'}$ are \emph{Morita equivalent}:
$A_r\otimes_\C\mathrm{Mat}_{r'}(\C)\simeq A_{r\cdot r'}$ exhibits each as a
matrix-algebra extension of the other.

\paragraph{Step (iii): adjoint Chern character.} For $\fgl_r$ acting on
itself by the adjoint, the degree-$3$ Chern character is
\[
\mathrm{ch}_3(\fgl_r,\mathrm{ad})=\tfrac{1}{6}\mathrm{tr}_{\mathrm{ad}}(F^3).
\]
Writing $F=F^a t_a$ with $t_a\in\fgl_r$ in the fundamental and using
$\mathrm{tr}_{\mathrm{ad}}(t_a t_b t_c)=(r-0)\mathrm{tr}_{\mathrm{fund}}(t_a t_b t_c)
+\text{mixed}$ (using $\mathrm{ad}_{\fgl_r}=\mathrm{fund}\otimes\mathrm{fund}^*
-\delta$ on trace-shifted part), one obtains the leading Zumino formula
\[
\mathrm{ch}_3(\fgl_r,\mathrm{ad})=\frac{r}{6}\mathrm{tr}_{\mathrm{fund}}(F^3)-\frac{1}{2}\mathrm{tr}_{\mathrm{fund}}(F^2)\mathrm{tr}_{\mathrm{fund}}(F).
\]
In particular, at $r=1$: $\mathrm{tr}_{\mathrm{fund}}(F)=F$ (scalar),
$\mathrm{tr}_{\mathrm{fund}}(F^3)=F^3$, and the two pieces cancel in
the $\chitop$-pairing, giving $\AAg_{\fgl_1}=0$. At $r\ge 2$, $\AAg$ is
non-zero and distinguishes $r$.

\paragraph{Step (iv): Costello's one-loop BV determinant.}
For $X=\C^3$ non-compact, the one-loop anomaly is computed via the
residue at $\infty$ (Costello 2013 \S4). On a compactification
$\overline X$ of $\C^3$ (e.g.\ $\overline X=\mathbb P^3$), the integral
\[
\log\det^{(1)}_\BV(\hCS_{\fg_r}(\overline X))
=\int_{\overline X}\mathrm{ch}_3(\fg_r)\wedge\mathrm{td}(\overline X)
\]
evaluates to a non-zero multiple of $r/6\cdot\chitop(\overline X)$ on
the adjoint-Chern-character term, plus the Zumino mixed term. Since
$\chitop(\mathbb P^3)=4$, the dominant piece is $(2/3)r$, growing with
$r$. Different $r$ give different BV determinants; the hCS theories are
genuinely distinct.

\paragraph{Step (v): no autoequivalence turns $V_r$ into $V_{r'}$.}
If $\phi\in\Aut(\Perf(\C^3))$, then $\phi(V_r)$ has endomorphism algebra
$R\End(\phi(V_r))\simeq\phi_*R\End(V_r)\simeq R\End(V_r)$ as dg algebras
(up to quasi-isomorphism induced by $\phi$). So $\phi(V_r)$ has the same
endomorphism algebra as $V_r$ up to isomorphism, which means
$\mathrm{Mat}_r(\C[x,y,z])$. This cannot equal $\mathrm{Mat}_{r'}(\C[x,y,z])$
for $r'\ne r$ (different degree-$0$ fibre dimensions). Therefore no
autoequivalence $\phi$ satisfies $\phi(V_r)\simeq V_{r'}$. The framings
$V_r$ for distinct $r$ are in distinct connected components of
$\Frame(\Perf(\C^3))$ that are \emph{not} related by $\Aut$-action;
they are related by Morita equivalence only at the level of categories
of modules, not at the level of generators.

\paragraph{Heal 5.} Inscribe the worked example as a theorem.

\begin{theorem}[$\C^3$ framings by rank; \ClaimStatusProvedHere]
\label{thm:C3-framings-rank}
Let $X=\C^3$, $\cC=\Perf(\C^3)$, and $V_r=\cO_{\C^3}^{\oplus r}$ for
$r\in\Z_{\ge 1}$. Then:
\begin{enumerate}[label=\textup{(\roman*)}]
\item Each $V_r$ is a compact generator of $\cC$, exhibiting the framing
$\cC\simeq\Perf(\mathrm{Mat}_r(\C[x,y,z]))$.
\item The endomorphism dg algebras $A_r=\mathrm{Mat}_r(\C[x,y,z])$ are
pairwise non-isomorphic but pairwise Morita-equivalent.
\item The framing Lie algebras $\fg_r=\fgl_r\otimes\C[x,y,z]$ are
pairwise non-isomorphic as Lie algebras.
\item The adjoint-Chern-character $\AAg_{\fgl_r}=\mathrm{ch}_3(\fgl_r,\mathrm{ad})$
is non-zero for $r\ge 2$ and grows with $r$.
\item Costello's one-loop BV determinant $\log\det^{(1)}_\BV(\hCS_{\fg_r}(\C^3))$
(computed on a compactification) grows linearly in $r$ in its leading
anomaly term, hence distinguishes the hCS theories.
\item No autoequivalence $\phi\in\Aut(\cC)$ satisfies
$\phi(V_r)\simeq V_{r'}$ for $r\ne r'$; the framings are in distinct
connected components of $\Frame(\cC)$ not linked by $\Aut$-action.
\end{enumerate}
Consequently the gauge datum $\fg_r$ is parametrised by a point of
$\Frame(\cC)\setminus\Frame(\cC)^\simeq$, \emph{not} by an element of
$\Aut(\cC)$.
\end{theorem}

\begin{proof}
(i) Compact generation and framing dg algebra: Steps~(i)--(ii) above.
(ii) Non-isomorphism: degree-$0$ fibre dimensions at the origin are
$r^2\ne(r')^2$. Morita equivalence: $(A_r)\otimes\mathrm{Mat}_{r'}=(A_{r\cdot r'})
\otimes\mathrm{Mat}_1=A_{r r'}$ on the nose.
(iii) Non-isomorphism of Lie algebras: $\dim_{\C[x,y,z]}\fg_r=r^2\ne(r')^2$.
(iv) Zumino formula: Step~(iii) above.
(v) BV determinant: Step~(iv) above, plus Costello 2013 Thm.~4.1.
(vi) Autoequivalence obstruction: Step~(v) above.
\end{proof}

\begin{corollary}[User's essay-claim, precise form]
\label{cor:users-essay-precise}
The essay-claim
\begin{quote}\emph{``Different $\fg$ on the same $X$ give different hCS
theories, different factorisation algebras, different BV determinants,
different one-loop anomalies; the gauge group lives in the 2-groupoid
of autoequivalences''}\end{quote}
is made precise by the substitution
\[
\boxed{\ \text{``the 2-groupoid of autoequivalences'' }\leadsto\
\text{``the $(\infty,1)$-category of framings'' }\ }
\]
viz., by Theorem~\ref{thm:framing-infty1}, the gauge datum
$(X,\fg)$ is indexed by a compact generator $V\in\Perf(X)$ with
$\fg=R\End^\bullet(V)_{\mathrm{Lie}}$, i.e.\ by an object of
$\Frame(\Perf(X))$. The core $\Frame(\Perf(X))^\simeq\simeq B\Aut(\Perf(X))$
is the 2-groupoid of autoequivalences, but this 2-groupoid \emph{does
not} index the gauge data; it acts on framings. Framings in distinct
connected components of $\Frame(\Perf(X))$ (e.g.\ $V_r$ for different
$r$) are not related by $\Aut$ and yield genuinely distinct hCS theories,
with BV determinants separated by Theorem~\ref{thm:invariant-tower}.
\end{corollary}

%=====================================================================
\subsection*{Synthesis: the healed formulation}
%=====================================================================

The precise, healed statement of the essay-claim, CG-voiced:

\begin{theorem}[Framing-indexed hCS; \ClaimStatusProvedHere, healed form
of user's essay]
\label{thm:framing-indexed-hCS}
Let $X$ be a smooth CY three-fold (toric or compact) and $\cC=\Perf(X)$
its bounded derived category of perfect complexes. Fix a framing
$V\in\cC$ (compact generator). The tilting algebra
$A_V=R\End^\bullet_\cC(V)$ is a dg algebra, and its associated Lie
algebra $\fg_V=A_{V,\mathrm{Lie}}$ is the gauge input to the Costello
holomorphic Chern--Simons theory $\hCS_{\fg_V}(X)$.

\begin{itemize}
\item[\textup{(1)}] \emph{Framing constellation.} The totality of
framings $V$ forms an $(\infty,1)$-category $\Frame(\cC)$, with
1-morphisms the Morita bimodules and 2-morphisms the bimodule maps.
Its core $\Frame(\cC)^\simeq=B\Aut(\cC)$ is the 2-groupoid of ambient
autoequivalences.
\item[\textup{(2)}] \emph{Invariant tower.} Bare-category invariants
($\chitop(X)$, $\mathrm{HH}^\bullet(\cC)$) are Morita-invariant; framing
invariants ($\rk\fg_V,\ \AAg_{\fg_V}=\mathrm{ch}_3(\fg_V,\mathrm{ad})$)
are not. The adjoint-Chern-character is the obstruction to
bare-category invariance of the gauge data.
\item[\textup{(3)}] \emph{One-loop BV anomaly.} The one-loop BV
determinant of $\hCS_{\fg_V}(X)$ is
\[
\log\det^{(1)}_\BV(\hCS_{\fg_V}(X))=\int_X\mathrm{ch}_3(\fg_V,\mathrm{ad})\wedge\mathrm{td}(X),
\]
which specialises on compact CY$_3$ to $\kanom(X,\fg_V)\propto\AAg(\fg_V)\chitop(X)$
(simplified form). This distinguishes framings in different connected
components of $\Frame(\cC)$.
\item[\textup{(4)}] \emph{Example $X=\C^3$.} Framings $V_r=\cO^{\oplus r}$
lie in pairwise-distinct $\Frame$-components not linked by $\Aut$; the
associated hCS theories $\hCS_{\fgl_r\otimes\C[x,y,z]}(\C^3)$ are
genuinely distinct.
\end{itemize}
The claim ``$\fg$ lives in $\Aut(\cC)$'' is false as literally stated;
the correct statement is ``$\fg$ is a framing in $\Frame(\cC)$, with
$\Aut(\cC)=\Frame(\cC)^\simeq$ acting but not indexing''.
\end{theorem}

\begin{corollary}[The CY-to-chiral functor sees framings, not only
$\cC$; \ClaimStatusSharpened]
\label{cor:phi-sees-framings}
In the notation of the manuscript main.tex, the CY-to-chiral functor
\[
\Phi:\mathrm{CY}\text{-cat}_d\longrightarrow\mathrm{ChirAlg}
\]
is not a functor of bare categories: its input is a \emph{framed CY
category} $(\cC,V)$, not just $\cC$. Different framings of the same
$\cC$ produce different chiral outputs, even at fixed $d$ and fixed $X$.
This resolves the apparent tension between
\begin{itemize}
\item[$\bullet$] ``$\Phi$ is a functor on CY categories'' (which is
true at the level of cores $\Frame^\simeq$, i.e.\ autoequivalences
permute framings), and
\item[$\bullet$] ``$\Phi$ output depends on $\fg$'' (which is true at
the level of framings themselves).
\end{itemize}
The correct domain for $\Phi$ is $\Frame(\mathrm{CY}\text{-cat}_d)$, the
$(\infty,1)$-category of framed CY categories. When the manuscript
writes $\Phi(X)$ without specifying a framing, the convention is
$V=\cO_X$ (trivial framing, giving $\fg_V=$ Hochschild Lie algebra of
functions on $X$).
\end{corollary}

\begin{remark}[Connection to AP-CY scope patterns]
Corollary~\ref{cor:phi-sees-framings} is the $\Frame$-level strengthening
of AP-CY discipline: Pattern~273 (``$\Phi$ functor vs object-level
correspondence'') in the CLAUDE.md is the \emph{object-level} (framing)
form; the functorial form lives on cores, i.e.\ on autoequivalences.
Neither is a hierarchy; they are two distinct statements about two
distinct $(\infty,1)$-categorical structures.
\end{remark}

\begin{remark}[Connection to wave12\_a3\_MO\_E2\_etingof.tex]
The MO $R$-matrix half-braiding identified there lives on
$\cZ(\Rep^{E_1}(Y(\widehat{\fgl}_1)))$, the derived centre of a module
category over a specific framed CY algebra. The framing is $V_r$ for
$Y(\widehat{\fgl}_r)$ (since
$Y(\widehat{\fgl}_r)$ is the CoHA of $\Perf(\C^3)$ with framing $V_r$).
Varying $r$ (the framing) varies the Yangian and its centre, as in
Theorem~\ref{thm:C3-framings-rank}.
\end{remark}

\begin{remark}[Connection to CoHA\_to\_W\_infty\_treatise.tex]
In the CoHA treatise, the CoHA construction on $\C^3$ with framing
$V=\cO^{\oplus r}$ gives the affine Yangian of $\fgl_r$. Varying $r$
gives different Yangians, as framings, all ``of the same bare category''.
\end{remark}

%=====================================================================
\subsection*{Summary (\(\le\)400 words)}
%=====================================================================

\paragraph{What was asserted.} Different gauge Lie algebras $\fg$ on the
same smooth CY three-fold $X$ give different holomorphic Chern--Simons
theories, different factorisation algebras, different BV determinants,
different one-loop anomalies $\kanom\propto\AAg\chitop(X)$ --- yet the
bare category $D^b\Coh(X)$ is independent of $\fg$. The essay proposed
``the 2-groupoid of autoequivalences'' as the locus of gauge data.

\paragraph{What survived five attack--heal cycles.}
\begin{enumerate}
\item The gauge datum is not in $\Aut(\cC)$. It is a \emph{framing}: a
compact generator $V\in\cC$ with $\fg=R\End^\bullet(V)_{\mathrm{Lie}}$
(Cycle~1).
\item The correct ambient structure is $\Frame(\cC)$, an
$(\infty,1)$-category, not a 2-groupoid. Its core $\Frame(\cC)^\simeq\simeq B\Aut(\cC)$
is the 2-groupoid of autoequivalences, but this core does not index
framings; it acts on them (Cycle~2).
\item ``External $=$ internal'' for gauge data holds iff Bondal--Van den
Bergh tilting applies: $\fg$ must arise as $R\End^\bullet(V)$ for some
compact generator. For $\fg=\fgl_r\otimes\cO_X$ this holds via
$V=\cO_X^{\oplus r}$. For abstract non-matrix-type $\fg$, it may fail
(Cycle~3).
\item The adjoint-Chern-character $\AAg=\mathrm{ch}_3(\fg,\mathrm{ad})$
is \emph{framing-dependent}, not Morita-invariant. The bare category
sees only Hochschild-level invariants; the framing sees $\AAg$; the hCS
theory sees the full anomaly integral (Cycle~4).
\item Worked example $X=\C^3$, $V_r=\cO^{\oplus r}$: endomorphism
algebras $\mathrm{Mat}_r(\C[x,y,z])$ are pairwise non-isomorphic but
Morita-equivalent; $\AAg_{\fgl_r}$ grows in $r$; BV determinants grow
linearly; no autoequivalence connects framings of different rank
(Cycle~5).
\end{enumerate}

\paragraph{Precise healed statement.} Theorem~\ref{thm:framing-indexed-hCS}.
The gauge input to hCS on a CY$_3$ is a framed CY category $(\cC,V)$,
not a bare $\cC$. The CY-to-chiral functor $\Phi$ takes framed CY
categories as input; different framings of the same $\cC$ give different
chiral outputs, and different Costello anomalies, all detectable by
$\AAg$ pairing against $\chitop$.

\paragraph{Voice.} The ``equals sign is a theorem''. The essay's
statement $\fg\in\Aut(\cC)$ is false; its mathematical content is
$\fg\in\mathrm{Ob}(\Frame(\cC))$, and this is a sharper, true theorem.
